
% -----------------------------------------------------------------------------
% Detailed charged-fermion QED scattering examples
% -----------------------------------------------------------------------------

一光子交换得到 Coulomb 势以后，还必须继续计算相对论可观测量。严格使用
\eqnrefs{eq:reduced-si-physical-p-r10,eq:gem-alpha-si-r10,eq:si-reduced-propagators-r10}
已经定义的 SI 约化图规则。四动量始终是物理动量，质量壳始终写成 $p^2=m^2c^2$；无量纲电磁耦合使用
\begin{equation}
 g_{\rm em}=\frac{e}{\sqrt{\varepsilon_0\hbar c}}=\sqrt{4\pi\alpha},
 \qquad
 \alpha=\frac{e^2}{4\pi\varepsilon_0\hbar c},
 \label{eq:alpha-si-qed-r10}
\end{equation}
其中 $e>0$ 的 SI 单位是 C，而电子实际电荷固定为 $\qe=-e$。图中使用的 $g_{\rm em}$ 是从 SI 电荷抽出的无量纲耦合；质量壳、传播子和截面则继续显式保留 $\hbar$ 与 $c$，因此每一步都可以直接核对量纲。

树级散射从已知的入射、出射粒子开始。它们的量子数决定哪些局域顶角允许出现，四动量流向则确定内部传播子的动量。由此写出 $\mathcal M$ 后，未观测自旋或极化必须在振幅平方层面平均与求和；只有得到 Lorentz 不变量 $\overline{|\mathcal M|^2}$，选择质心系或实验室系才只是运动学化简。最后与质量壳相空间和入射通量组合，才得到可以直接与实验比较的截面。

先取最简单的可区分费米子过程
$e^-(p_1)+\mu^-(p_2)\to e^-(p_3)+\mu^-(p_4)$。最低阶只有 $t$ 通道单光子交换，其中 $t=(p_1-p_3)^2$。对未偏振束流，初态电子和 $\mu$ 子各平均两种自旋，末态自旋求和再用 Dirac 完备关系闭合成两条迹，因此振幅平方的通用起点是
\begin{align}
 \overline{|\mathcal M_t|^2}
 =&\frac{g_{\rm em}^4}{t^2}\frac14
 \Tr[(\slashed p_3+\me c)\gamma^\mu(\slashed p_1+\me c)\gamma^\nu]\nonumber\\
 &\hspace{2.4cm}\times
 \Tr[(\slashed p_4+m_\mu c)\gamma_\mu(\slashed p_2+m_\mu c)\gamma_\nu].
 \label{eq:emu-M2-trace}
\end{align}
把两条 Dirac 迹缩并并用 Mandelstam 变量消去标量积，可先得到完整质量依赖的 Lorentz 不变量结果
\begin{equation}
 \boxed{
 \overline{|\mathcal M_t|^2}
 =\frac{2g_{\rm em}^4}{t^2}
 \left\{
 \left[s-(\me^2+m_\mu^2)c^2\right]^2
 +\left[u-(\me^2+m_\mu^2)c^2\right]^2
 +2t(\me^2+m_\mu^2)c^2
 \right\}.}
 \label{eq:emu-M2-massive-invariants-dp49}
\end{equation}
这个式子仍未选择参考系，也没有忽略电子或 $\mu$ 子质量。它把长程传播子产生的 $1/t^2$ 与外部费米子质量、自旋结构明确分开。

只有在
\begin{equation*}
 s\gg \me^2c^2,\qquad s\gg m_\mu^2c^2
\end{equation*}
的超相对论质心系中，才可把外腿质量相对于 $s$ 忽略，并使用
$t=-s(1-\cos\theta)/2$、$u=-s(1+\cos\theta)/2$。二体相空间母式随即给出
\begin{equation}
 \boxed{
 \frac{\dd\sigma(e\mu)}{\dd\Omega}
 =\frac{\hbar^2\alpha^2}{2s}
 \frac{4+(1+\cos\theta)^2}{(1-\cos\theta)^2}.}
 \label{eq:emu-angular-massless}
\end{equation}
分母中的 $(1-\cos\theta)^2$ 直接来自 $t^{-2}$，所以 $\theta\to0$ 的前向增强是无质量光子长程传播的运动学表现；有限束流角接受度、屏蔽或有限波包宽度会在实际测量中提供前向调节尺度。

\begin{figure}[htbp]
 \centering
 \begin{tikzpicture}[scale=0.92, baseline=(current bounding box.center)]
\coordinate (a) at (-3,1.5); \coordinate (b) at (-1,0.4); \coordinate (c) at (-3,-1.5); \coordinate (d) at (-1,-0.4);
\coordinate (e) at (1,0.4); \coordinate (f) at (3,1.5); \coordinate (g) at (1,-0.4); \coordinate (h) at (3,-1.5);
\draw[fermion] (a)--(b); \draw[fermion] (b)--(f);
\draw[fermion] (c)--(d); \draw[fermion] (d)--(h);
\draw[photon] (b)--(d);
\node[left] at (a) {$e^-(p_1)$}; \node[left] at (c) {$\mu^-(p_2)$};
\node[right] at (f) {$e^-(p_3)$}; \node[right] at (h) {$\mu^-(p_4)$};
\node[right=2pt] at (-1,-0.0) {$q^2=t$};
\end{tikzpicture}

 \caption{$e^-\mu^-\to e^-\mu^-$ 的最低阶只有 $t$ 通道光子交换。可区分外腿使它成为学习电流缩并、自旋平均与 $1/t^2$ 前向峰最直接的 QED 例子。}
 \label{fig:emu-scattering}
\end{figure}

\begin{figure}[htbp]
 \centering
 \includegraphics[width=.76\linewidth]{emu_angular_distribution.png}
 \caption{超相对论 $e^-\mu^-$ 树级角分布，使用 \eqnrefs{eq:emu-angular-massless} 计算并以 $\theta=90^\circ$ 归一化。小角的强增长来自 $t$ 通道无质量光子极点。}
 \label{fig:emu-angular}
\end{figure}

接着把初态改为粒子--反粒子，考虑
$e^-(p_1)+e^+(p_2)\to\mu^-(p_3)+\mu^+(p_4)$。此时最低阶允许 $s$ 通道湮灭，内部光子携带总四动量 $p_1+p_2$，故分母为 $s=(p_1+p_2)^2$。沿费米子数箭头读取外部旋量，约化振幅为
\begin{equation}
 \boxed{
 \mathcal M_s
 =\frac{g_{\rm em}^2}{s}
 [\bar v(p_2)\gamma^\mu u(p_1)]
 [\bar u(p_3)\gamma_\mu v(p_4)].}
 \label{eq:ee-mumu-M}
\end{equation}

这个过程把振幅、自旋动力学、相空间和角积分四个层次清楚地串在一起。真实 $\mu^+\mu^-$ 末态要求 $s\ge4m_\mu^2c^2$；在整个阈值区域都有 $s\gg\me^2c^2$，因此可先忽略初态电子质量，但完整保留末态 $\mu$ 子质量。先用一个无量纲速度参数记录末态相空间的开启程度，定义
\begin{equation}
 \beta_{\rm rel}^{(\mu)}
 \equiv\sqrt{1-\frac{4m_\mu^2c^2}{s}},
 \qquad
 |\mathbf p_f|=\frac{\sqrt s}{2}\beta_{\rm rel}^{(\mu)},
 \qquad
 |\mathbf p_i|=\frac{\sqrt s}{2},
 \label{eq:mumu-beta-si-r10}
\end{equation}
则 $c\beta_{\rm rel}^{(\mu)}$ 是质心系末态 $\mu$ 子速率。初态自旋平均和末态自旋求和只处理两条 Dirac 电流的内部自由度，因此先得到与相空间无关的 $\overline{|\mathcal M|^2}$：
\begin{equation}
 \boxed{
 \overline{|\mathcal M|^2}
 =g_{\rm em}^4\left[
 1+\cos^2\theta
 +(1-(\beta_{\rm rel}^{(\mu)})^2)\sin^2\theta
 \right].}
 \label{eq:mumu-M2-massive-angle-si-r10}
\end{equation}
再乘二体相空间与初末态质心动量之比，才得到微分截面：
\begin{equation}
 \boxed{
 \frac{\dd\sigma}{\dd\Omega}
 =\frac{\hbar^2\alpha^2\beta_{\rm rel}^{(\mu)}}{4s}
 \left[
 1+\cos^2\theta
 +(1-(\beta_{\rm rel}^{(\mu)})^2)\sin^2\theta
 \right],}
 \label{eq:mumu-dsigma-massive-si-r10}
\end{equation}
以及
\begin{equation}
 \boxed{
 \sigma
 =\frac{4\pi\hbar^2\alpha^2}{3s}
 \beta_{\rm rel}^{(\mu)}\left(1+\frac{2m_\mu^2c^2}{s}\right).}
 \label{eq:mumu-total-massive-si-r10}
\end{equation}
阈值 $\beta_{\rm rel}^{(\mu)}\to0$ 时，总截面由末态相空间线性关闭，而角分布趋于各向同性；阈值形状因而首先反映可用二体相空间。再取 $m_\mu^2c^2/s\to0$ 才得到超相对论极限
\begin{equation}
 \boxed{
 \sigma(e^+e^-\to\mu^+\mu^-)
 =\frac{4\pi\hbar^2\alpha^2}{3s},}
 \label{eq:ee-mumu-total}
\end{equation}
且微分角分布正比于 $1+\cos^2\theta$。与 $t$ 通道散射不同，该 $s$ 通道振幅没有前向极点；固定束流能量时，$s$ 不随散射角变化，角结构来自两条矢量流的自旋代数。

\begin{figure}[htbp]
 \centering
 \begin{tikzpicture}[scale=0.92,baseline=(current bounding box.center),
 fermion/.style={thick,postaction={decorate},decoration={markings,mark=at position .56 with {\arrow{Latex}}}},
 antifermion/.style={thick,postaction={decorate},decoration={markings,mark=at position .44 with {\arrowreversed{Latex}}}},
 photon/.style={thick,decorate,decoration={snake,amplitude=1.5pt,segment length=7pt}}]
\coordinate (L) at (-1.2,0); \coordinate (R) at (1.2,0);
\draw[fermion] (-3,1.5)--(L); \draw[antifermion] (-3,-1.5)--(L);
\draw[photon] (L)--(R);
\draw[fermion] (R)--(3,1.5); \draw[antifermion] (R)--(3,-1.5);
\node[left] at (-3,1.5) {$e^-$}; \node[left] at (-3,-1.5) {$e^+$};
\node[right] at (3,1.5) {$\mu^-$}; \node[right] at (3,-1.5) {$\mu^+$};
\node[above] at (0,0.15) {$q^2=s$};
\end{tikzpicture}

 \caption{$e^-e^+\to\mu^-\mu^+$ 的 $s$ 通道光子湮灭。阈值由真实 $\mu^+\mu^-$ 二体相空间决定。}
 \label{fig:ee-mumu}
\end{figure}

\begin{figure}[htbp]
 \centering
 \includegraphics[width=.72\linewidth]{ee_mumu_angular_distribution.png}
 \caption{$e^-e^+\to\mu^-\mu^+$ 在 $s\gg m_\mu^2c^2$ 下的归一化角分布 $1+\cos^2\theta$。有限质量时，\eqnrefs{eq:mumu-dsigma-massive-si-r10} 还包含阈值附近的各向同性分量。}
 \label{fig:ee-mumu-angular}
\end{figure}
\begin{exbox}[$10\,\mathrm{GeV}$ 的 $e^+e^-\to\mu^+\mu^-$]
高能极限\eqnrefs{eq:ee-mumu-total} 用上述物理四动量变量写成
\begin{equation*}
 \sigma=\frac{4\pi\hbar^2\alpha^2}{3s},
 \qquad
 s=\left(\frac{E_{\rm CM}}{c}\right)^2,
\end{equation*}
所以也可以直接改写为能量单位更方便的形式
\begin{equation}
 \boxed{
 \sigma=\frac{4\pi\alpha^2(\hbar c)^2}{3E_{\rm CM}^2}.}
 \label{eq:ee-mumu-total-energy-form-dp61}
\end{equation}
取 $E_{\rm CM}=10\,\mathrm{GeV}$、$\alpha^{-1}=137.036$ 与
$\hbar c=0.19732698\,\mathrm{GeV\,fm}$，逐项代入：
\begin{align*}
 \sigma
 &=\frac{4\pi}{3}
 \left(\frac1{137.036}\right)^2
 \frac{(0.19732698\,\mathrm{GeV\,fm})^2}
 {(10\,\mathrm{GeV})^2},\\
 &=8.69\times10^{-8}\,\mathrm{fm^2},\\
 &=0.869\,\mathrm{nb}
 =8.69\times10^2\,\mathrm{pb}.
\end{align*}
这个数量级可以直接用来检查程序中的相空间归一化：若结果少了 $(\hbar c)^2$，截面甚至不会具有面积量纲；若误把 $E_{\rm CM}$ 而不是 $E_{\rm CM}^2$ 放入分母，高能标度也会立即错误。

\figref{fig:ee-mumu-total-energy-dp61} 把有限 $m_\mu$ 的树级结果与 $m_\mu\to0$ 的 $1/E_{\rm CM}^2$ 极限放在同一图中。这里使用低能精细结构常数并只保留单光子树图；真空极化造成的运行耦合和电弱 $Z$ 交换属于后续更高精度修正。
\end{exbox}

\begin{figure}[htbp]
 \centering
 \includegraphics[width=.78\linewidth]{ee_mumu_total_cross_section_dp61.png}
 \caption{$e^+e^-\to\mu^+\mu^-$ 的树级总截面。红线保留 $m_\mu$，因此在 $E_{\rm CM}=2m_\mu c^2$ 阈值处由二体相空间关闭；虚线为 $m_\mu\to0$ 后的 $1/E_{\rm CM}^2$ 渐近式。图只显示纯 QED 树级结构。}
 \label{fig:ee-mumu-total-energy-dp61}
\end{figure}



若初末态都是电子，末态标签不可区分。最低阶同时存在 $t$、$u$ 两种外腿连接，并且交换两个外部费米子带来相对负号，所以必须先在振幅层计算 $\mathcal M_t-\mathcal M_u$，再平方得到干涉项。在 $s\gg \me^2c^2$ 的超相对论极限，并对不可区分末态采用单次计数的角域约定，得到
\begin{equation}
 \boxed{
 \frac{\dd\sigma_{\mathrm{Moller}}}{\dd\Omega}
 =\frac{\hbar^2\alpha^2}{2s}
 \frac{(3+\cos^2\theta)^2}{\sin^4\theta}.}
 \label{eq:moller-angular}
\end{equation}
$\theta\to0$ 与 $\theta\to\pi$ 的两端增强分别继承自 $t$、$u$ 通道；分子中的 $(3+\cos^2\theta)^2$ 则包含两条不可区分路径的费米交换干涉。

\begin{figure}[htbp]
 \centering
 \begin{tikzpicture}[scale=.78, baseline=(current bounding box.center)]
% t channel
\begin{scope}[xshift=-4.0cm]
\draw[fermion] (-2,1.5)--(-.6,.5)--(2,1.5);
\draw[fermion] (-2,-1.5)--(-.6,-.5)--(2,-1.5);
\draw[photon] (-.6,.5)--(-.6,-.5);
\node[above] at (-.6,.65) {$t$};
\end{scope}
\node at (0,0) {$-$};
% u channel: identical topology with exchanged external labels.  Keeping the
% fermion strokes non-crossing makes the graph topology explicit; the exchange
% is carried by the external-state assignment in the accompanying amplitude.
\begin{scope}[xshift=4.0cm]
\draw[fermion] (-2,1.5)--(-.6,.5)--(2,1.5);
\draw[fermion] (-2,-1.5)--(-.6,-.5)--(2,-1.5);
\draw[photon] (-.6,.5)--(-.6,-.5);
\node[above] at (-.6,.65) {$u$};
\end{scope}
\end{tikzpicture}

 \caption{M{\o}ller 散射的 $t$ 与 $u$ 两张树图。相对负号来自相同费米子交换，而不是光子传播子。}
 \label{fig:moller-two}
\end{figure}

\begin{figure}[htbp]
 \centering
 \includegraphics[width=.74\linewidth]{moller_angular_distribution.png}
 \caption{超相对论 M{\o}ller 散射的角分布，由 \eqnrefs{eq:moller-angular} 计算并以 $90^\circ$ 归一化。$\theta\to0$ 与 $\theta\to\pi$ 的双侧增强分别对应 $t$、$u$ 两个无质量光子交换极点。}
 \label{fig:moller-angular}
\end{figure}

Bhabha 散射 $e^-(p_1)e^+(p_2)\to e^-(p_3)e^+(p_4)$ 与 M{\o}ller 散射的关键区别是末态电子和正电子可区分。最低阶仍有两张图，但现在分别是 $s$ 通道湮灭和 $t$ 通道交换。按照同一套外腿与费米子流约定，两条图必须先在振幅层相加：
\begin{equation}
 \boxed{
 \mathcal M_{\rm Bhabha}
 =\frac{g_{\rm em}^2}{s}
 [\bar v_2\gamma^\mu u_1][\bar u_3\gamma_\mu v_4]
 -\frac{g_{\rm em}^2}{t}
 [\bar u_3\gamma^\mu u_1][\bar v_2\gamma_\mu v_4].}
 \label{eq:bhabha-tree-amplitude-r68}
\end{equation}
在 $s\gg \me^2c^2$ 的超相对论极限，对初态自旋平均、末态自旋求和并保留 $s$--$t$ 干涉后得到
\begin{equation}
 \boxed{
 \frac{\dd\sigma_{\rm Bhabha}}{\dd\Omega}
 =\frac{\hbar^2\alpha^2}{4s}
 \frac{(3+\cos^2\theta)^2}{(1-\cos\theta)^2}.}
 \label{eq:bhabha-angular-r68}
\end{equation}
因此前向峰来自 $t\to0$ 的无质量光子交换，而较大角度处的形状同时包含湮灭通道及其与 $t$ 通道的干涉。与 M{\o}ller 情形不同，这里不存在全同末态的 $1/2!$ 重复计数因子。

\begin{figure}[htbp]
 \centering
 \begin{tikzpicture}[scale=.78,baseline=(current bounding box.center),
 fermion/.style={thick,postaction={decorate},decoration={markings,mark=at position .56 with {\arrow{Latex}}}},
 antifermion/.style={thick,postaction={decorate},decoration={markings,mark=at position .44 with {\arrowreversed{Latex}}}},
 photon/.style={thick,decorate,decoration={snake,amplitude=1.4pt,segment length=7pt}}]
% s
\begin{scope}[xshift=-4.0cm]
\coordinate (L) at (-.7,0);\coordinate (R) at (.7,0);
\draw[fermion] (-2,1.4)--(L);\draw[antifermion] (-2,-1.4)--(L);\draw[photon] (L)--(R);\draw[fermion] (R)--(2,1.4);\draw[antifermion] (R)--(2,-1.4);\node[above] at (0,.15){$s$};
\end{scope}
\node at (0,0) {$+$};
% t
\begin{scope}[xshift=4.0cm]
\draw[fermion] (-2,1.4)--(-.6,.45)--(2,1.4);\draw[antifermion] (-2,-1.4)--(-.6,-.45)--(2,-1.4);\draw[photon] (-.6,.45)--(-.6,-.45);\node[above] at (-.6,.6){$t$};
\end{scope}
\end{tikzpicture}

 \caption{Bhabha 散射的 $s$ 通道湮灭与 $t$ 通道交换。}
 \label{fig:bhabha-two}
\end{figure}

\begin{figure}[htbp]
 \centering
 \includegraphics[width=.74\linewidth]{bhabha_angular_distribution.png}
 \caption{超相对论 Bhabha 树级角分布，以 $90^\circ$ 为归一化基准。前向峰来自 $t$ 通道光子交换；与 M{\o}ller 散射的前后向对称增强不同，末态 $e^-$ 与 $e^+$ 可区分。}
 \label{fig:bhabha-angular}
\end{figure}

这些例子已经说明为什么只背一张“QED 顶角 $-\ii e\gamma^\mu$”远远不够：同一个顶角在 $e\mu$、湮灭、M{\o}ller、Bhabha 四个过程中分别遇到单一交换、单一湮灭、费米子交换反对称性、以及 $s$--$t$ 干涉。Feynman 图真正的计算能力来自“图拓扑 + 外态统计 + 自旋/偏振代数 + 相空间”四者同时使用。

